\documentclass[12pt,a4paper,oneside]{report}
\usepackage[utf8]{inputenc}
\usepackage[T1]{fontenc}
\usepackage[spanish,es-nodecimaldot,es-noquoting]{babel}
\usepackage{lmodern,microtype,amsmath,amssymb,mathtools}
\usepackage[margin=2.2cm,headheight=16pt]{geometry}
\usepackage{parskip,fancyhdr,booktabs,tabularx,longtable,array,enumitem}
\usepackage{xcolor,tikz,listings,tcolorbox,titlesec,needspace}
\tcbuselibrary{breakable,skins}
\usetikzlibrary{arrows.meta,positioning,calc,fit,matrix,decorations.pathreplacing}
\usepackage[hidelinks,unicode]{hyperref}
\usepackage{bookmark}
\makeatletter
\renewcommand{\@pnumwidth}{2.3em}
\renewcommand{\@tocrmarg}{3.3em}
\renewcommand{\l@chapter}[2]{\addpenalty{-\@highpenalty}\vskip1em\@plus\p@\@dottedtocline{0}{0pt}{2.5em}{\bfseries #1}{\bfseries #2}}
\renewcommand{\l@section}{\@dottedtocline{1}{1.5em}{3.3em}}
\makeatother
\pdfstringdefDisableCommands{\def\times{x}\def\ensuremath#1{#1}}
\hypersetup{pdftitle={Compiladores para aprendizaje automático: desde cero},pdfauthor={Material didáctico independiente},pdfsubject={Curso completo con compilador Lumbre, CPU, GPU, autodiferenciación y modelos de lenguaje}}
\definecolor{azul}{RGB}{25,45,133}
\definecolor{verde}{RGB}{16,108,53}
\definecolor{naranja}{RGB}{174,90,9}
\definecolor{violeta}{RGB}{94,35,125}
\definecolor{rojo}{RGB}{166,30,35}
\definecolor{petroleo}{RGB}{18,103,106}
\newcommand{\cajadef}[2]{\newtcolorbox{#1}[1][]{enhanced,breakable,colback=#2!4!white,colframe=#2,fonttitle=\bfseries,title={##1},coltitle=white,colbacktitle=#2,boxrule=.8pt,arc=2mm,left=7pt,right=7pt,top=5pt,bottom=5pt,before skip=9pt,after skip=9pt}}
\cajadef{idea}{azul}
\cajadef{definicion}{azul}
\cajadef{traduccion}{naranja}
\cajadef{ejemplo}{verde}
\cajadef{practica}{violeta}
\cajadef{cuidado}{rojo}
\cajadef{comprueba}{petroleo}
\cajadef{nota}{gray}
\pagestyle{fancy}\fancyhf{}
\fancyhead[L]{\small\itshape Compiladores de ML: desde cero}
\fancyhead[R]{\small\itshape Edición 05-10-2026}
\fancyfoot[C]{\thepage}
\renewcommand{\headrulewidth}{.4pt}
\fancypagestyle{plain}{\fancyhf{}\fancyfoot[C]{\thepage}\renewcommand{\headrulewidth}{0pt}}
\titleformat{\chapter}[display]{\normalfont\huge\bfseries}{\color{azul}\large CAPÍTULO \thechapter}{8pt}{}
\titlespacing*{\chapter}{0pt}{0pt}{16pt}
\titleformat{\section}{\normalfont\Large\bfseries}{\thesection}{.7em}{}
\setcounter{tocdepth}{1}\setcounter{secnumdepth}{2}
\setlength{\emergencystretch}{3em}
\setlength{\LTpre}{8pt}\setlength{\LTpost}{8pt}
\renewcommand{\arraystretch}{1.16}
\newcolumntype{Y}{>{\raggedright\arraybackslash}X}
\newcolumntype{P}[1]{>{\raggedright\arraybackslash}p{#1}}
\setlist{leftmargin=*,topsep=5pt,itemsep=4pt}
\linespread{1.045}
\lstset{basicstyle=\ttfamily\footnotesize,breaklines=true,breakatwhitespace=false,columns=fullflexible,keepspaces=true,showstringspaces=false,frame=single,rulecolor=\color{gray!40},backgroundcolor=\color{gray!4},tabsize=4,xleftmargin=3pt,xrightmargin=3pt,aboveskip=9pt,belowskip=9pt,upquote=true,postbreak=\mbox{\textcolor{gray}{$\hookrightarrow$}\space},literate={á}{{\'a}}1 {é}{{\'e}}1 {í}{{\'i}}1 {ó}{{\'o}}1 {ú}{{\'u}}1 {ñ}{{\~n}}1 {ü}{{\"u}}1 {Á}{{\'A}}1 {É}{{\'E}}1 {Í}{{\'I}}1 {Ó}{{\'O}}1 {Ú}{{\'U}}1 {Ñ}{{\~N}}1 {¿}{{?`}}1 {¡}{{!`}}1}
\tikzset{>=Stealth,box/.style={draw=azul,fill=azul!4,rounded corners=2pt,align=center,text width=3.9cm,minimum height=1.05cm,font=\small},flow/.style={->,thick,draw=azul},dot/.style={circle,draw=azul,fill=azul!4,minimum size=9mm,font=\small}}
\newcommand{\fuentes}[1]{\par\smallskip{\footnotesize\color{gray!75!black}\textbf{Fuentes y alcance.} #1\par}}
\newcommand{\mat}[1]{\mathbf{#1}}
\newcommand{\R}{\mathbb{R}}
\newcommand{\N}{\mathbb{N}}
\newcommand{\dd}{\mathrm{d}}
\begin{document}
\hypersetup{pageanchor=false}
\begin{titlepage}
\centering\vspace*{.8cm}
{\large INGENIERÍA INFORMÁTICA · SISTEMAS DE INTELIGENCIA ARTIFICIAL\par}
\vspace{1.0cm}
{\Huge\bfseries Compiladores para\\[.3cm]aprendizaje automático\par}
\vspace{.8cm}
{\LARGE\color{azul}\bfseries Desde cero\par}
\vspace{.7cm}
{\Large Comprender una operación. Construir su compilador.\\[.2cm]Entrenar un modelo y explicar cada resultado.\par}
\vspace{1cm}\rule{.88\textwidth}{.8pt}\vspace{.7cm}
\begin{minipage}{.85\textwidth}\large
De una suma en C a un Transformer: representaciones intermedias, reescrituras, memoria, kernels CPU y GPU, autodiferenciación y entrenamiento.
\par\medskip
Con el compilador \textbf{Lumbre}, laboratorios resueltos, proyectos de investigación, código completo y un mapa crítico de tinygrad, NVIDIA, AMD, Linux y Huawei Ascend.
\end{minipage}
\vspace{.65cm}\rule{.88\textwidth}{.8pt}
\vfill
{\large Curso de un semestre · Itinerario preparatorio incluido\par}
\vspace{.3cm}{\small Edición ampliada con el syllabus: 5 de octubre de 2026\par}
\vspace{.15cm}{\small Corte de investigación original: 1 de octubre de 2026\par}
\vspace{.35cm}
\begin{minipage}{.87\textwidth}\footnotesize
Material didáctico independiente. Los ejemplos, diagramas y el compilador Lumbre son elaboración original. Incluye validación en CPU y CUDA sobre una RTX 5070 Ti Laptop, comparaciones reproducibles y entrenamiento de un decoder pequeño. Los resultados externos se atribuyen a sus fuentes. HIP queda pendiente de validación en hardware compatible.
\end{minipage}
\end{titlepage}
\pagenumbering{arabic}
\hypersetup{pageanchor=true}
